

\chapter{Applications of martingale theory}

\section{Pólya's urn}
We now revisit the P\'olya urn experiment discussed
in Chapter 1 in which we start with $a$ white balls
and $b$ black balls, and repeatedly sample balls at
random from the urn, at each step putting back the
ball we sampled and adding another ball with the same
color. We showed in that discussion that the
proportion

\[
M_n = \frac{X_n}{n+a+b}
\]

\noindent of white balls in the urn is a martingale, which
takes values in $[0,1]$. By the martingale
convergence theorem we see that the limit
$M = \lim_{n\to\infty}$ exists almost surely. By an
explicit computation mentioned in Chapter 1, it
follows that $M$ is distributed according to a beta
distribution $\operatorname{Beta}(a,b)$.

\section{Recurrence of simple random walk on $\mathbb{Z}$}
Let $S_n = \sum_{k=1}^n X_k$ denote the simple random
walk on $\Z$; i.e., $S_0=0$ and $(X_n)_{n=1}^\infty$
are i.i.d.\ with
$\mathrm{P}(X_n=-1)=\mathrm{P}(X_n=1)=\tfrac12$.

\begin{thm}
The random walk is $\textbf{recurrent}$. That is, for
any $m\in\Z$ we have

\[
\mathrm{P}(S_n = m \textrm{ i.o.}) = 1.
\]

\end{thm}

\begin{proof}
Define a random variable
$N = \inf \{ n>0\,:\, S_n = -1 \}$ (on the event that
$S_n$ never visits $-1$, set $N=\infty$). Let
$M_n = S_{N\wedge n}$. By Theorem
\ref{thm-stopping-mart}, the sequence $(M_n)_{n=1}^\infty$
is a martingale with respect to the filtration
$\mathcal{G}_n = \sigma(X_1,\ldots,X_n)$.
Furthermore, $M_n \ge -1$ for all $n$, so by
Corollary \ref{mart-conv-cor}, the limit
$M=\lim_{n\to\infty} M_n$ exists almost surely. But
then we must have $M\equiv -1$ almost surely, since
any other limit is impossible (until the sequence
$S_{N\wedge n}$ stops at $-1$, it fluctuates by
$\pm 1$ at each step). This implies that $N<\infty$
almost surely, so the random walk is guaranteed to
visit $-1$ with probability $1$.
\end{proof}

\begin{xexercise}
Complete the proof by explaining why if the random
walk almost surely visits $-1$ when starting from $0$
then it almost surely visits every lattice point
$m\in\Z$, and therefore also almost surely visits
every $m\in\Z$ infinitely often.
\end{xexercise}

\noindent Note that the almost sure martingale limit
$M=\lim_{n\to\infty} S_{N\wedge n}$ in the proof
above satisfies $\mathbf{E} M=-1$, whereas
$\mathbf{E}(S_{N \wedge n})=0$ for all $n$. This is
therefore an example of a martingale that converges
almost surely, but not in $L_1$.

\section{Branching processes and the Galton-Watson tree}
The $\textbf{Galton-Watson tree}$, or
$\textbf{Galton-Watson process}$, (named after the
19th century statistics pioneers Francis Galton and
Henry William Watson) is a mathematical model for a
genealogical tree, or for the population statistics of
a unisexual animal species. It is the simplest in a
family of models known as
$\textbf{branching processes}$. The model was
originally conceived at a time when family names were
only passed on by male descendants and hence takes
only male offspring into account, leading to a
somewhat simpler (though socially anachronistic)
mathematical model than some later more realistic
variants.  (The original model is still used however
to model various biological and physical phenomena.)

Let $p_0, p_1, \ldots$ be nonnegative numbers such
that $\sum_k p_k = 1$, and let $X$ be a random
variable with $\mathrm{P}(X=k)=p_k$, $n=0,1,2,\ldots$.
We think of $p_k$ as the probability for a specimen
to have $k$ offspring, and refer to the distribution
of $X$ as the $\textbf{offspring distribution}$. Let
$(X_{n,m})_{n,m\ge 0}$ be a family of i.i.d.\ copies
of $X$. The Galton-Watson process is a sequence of
random variables $Z_0, Z_1, Z_2, \ldots$, where for
each $n$, $Z_n$ represents the number of
$n$th-generation descendants of the original species
patriarch at time $0$. $Z_n$ is defined by the
initial condition $Z_0=1$ together with the
recurrence relation

\[
Z_n = \sum_{m=1}^{Z_{n-1}} X_{n-1,m},
\]

\noindent corresponding to the assumption that each of the
$Z_{n-1}$$(n-1)$th-generation specimens bear a
number of offspring with the distribution of $X$,
with all offspring numbers being independent of each
other.

Let $\mu=\mathbf{E} X = \sum_k k p_k$ be the expected
number of offspring. We will prove:

\begin{thm}
\label{thm-galton-watson}
Let $E$ be the ``extinction event''

\[
E = \{ Z_n=0 \textnormal{ for all large enough }n\} = \{ \textnormal{the species eventually becomes extinct} \}.
\]

\noindent Then

\begin{enumerate}[label=\arabic*.]
\item If $\mu< 1$ then $\mathrm{P}(E) = 1$.

\item If $\mu=1$ and $\mathrm{P}(X=1)<1$ then
$\mathrm{P}(E)=1$.

\item If $\mu>1$ then $\mathrm{P}(E)<1$.
\end{enumerate}

\end{thm}

\noindent The reason martingale theory is relevant to the
problem is the following simple lemma.

\begin{lem}
The sequence $(Z_n/\mu^n)_{n=0}^\infty$ is a
martingale with respect to the filtration
$\mathcal{G}_n = \sigma(X_{\ell,m}\,:\, 0\le \ell< n,
m\ge 0)$.
\end{lem}

\begin{proof}
Compute
$\mathbf{E}\left(Z_{n+1}\,|\,\mathcal{G}_n\right)$,
using the standard properties of conditional
expectation:

\begin{align*}
\mathbf{E}\left(Z_{n+1}\,|\,\mathcal{G}_n\right) &= \mathbf{E}\left(\sum_{m=1}^{Z_{n}} X_{n,m}\,|\,\mathcal{G}_n\right)
=
\mathbf{E}\left( \sum_{k=0}^\infty  \mathbf{1}_{\{Z_n=k\}} \sum_{m=1}^{Z_{n}} X_{n,m}\,|\,\mathcal{G}_n\right)
\\
&= \sum_{k=0}^\infty
\mathbf{E}\left( \mathbf{1}_{\{Z_n=k\}} \sum_{m=1}^{k} X_{n,m}\,|\,\mathcal{G}_n\right)
= \sum_{k=0}^\infty \mathbf{1}_{\{Z_n=k\}}
\mathbf{E}\left(  \sum_{m=1}^{k} X_{n,m}\,|\,\mathcal{G}_n\right) \\
&=
\sum_{k=0}^\infty \mathbf{1}_{\{Z_n=k\}} k \mathbf{E}(X_{n,1}) = \mu \sum_{k=0}^\infty k \mathbf{1}_{\{Z_n=k\}} = \mu Z_n.
\end{align*}

\noindent Dividing by $\mu^{n+1}$ gives the result.
\end{proof}

\begin{proof}[Proof of parts 1 and 2 of Theorem \ref{thm-galton-watson}]
By the lemma we have $\mathbf{E}(Z_n)=\mu^n$. If
$\mu<1$ then

\[
\mathrm{P}(Z_n>0)=\mathrm{P}(Z_n\ge 1)\le \mathbf{E}(Z_n) = \mu^n \to 0.
\]

\noindent Since $E^c = \cap_{n=1}^\infty \{ Z_n>0 \}$, an
intersection of a decreasing sequence of sets, we
therefore get that
$\mathrm{P}(E^c)=\displaystyle\lim_{n\to\infty}
\mathrm{P}(Z_n>0)=0$,
which proves part 1.

To prove part 2, let

\[
Z_\infty = \lim_{n\to\infty} \mu^{-n} Z_n
\]

\noindent be the limit of the nonnegative martingale
$\mu^{-n} Z_n$, guaranteed to exist almost surely by
Corollary \ref{mart-conv-cor}. In the case when $\mu=1$,
$Z_\infty$ is the limit of the sequence $Z_n$ itself.
Under the assumption that $\mathrm{P}(X=1)<1$, this
limit can only be $0$, since on the event that
$\{Z_n=k>0\}$, $|Z_{n+1}-k|\ge 1$ with positive
probability.
\end{proof}

\noindent To prove the final claim regarding the behavior when
$\mu>1$, associate with the random variable $X$
(which takes nonnegative integer values) a function

\[
\varphi(z) = \mathbf{E}(z^X) = \sum_{k=0}^\infty p_k z^k, \qquad (0\le z\le 1).
\]

\noindent The function $\varphi(z)$ is called the \textbf{generating
function of $X$}. It turns out that the way it
encodes information about the distribution of $X$ is
especially suited to the problem at hand, because of
the following lemma.

\begin{lem}
The probability for the Galton-Watson process to
become extinct by the $n$th generation is given by

\[
\mathrm{P}(Z_n = 0) = \varphi^{(n)}(0) = (\overbrace{\varphi\circ \varphi\circ \ldots \circ \varphi}^{\textrm{$n$ times}})(0),
\]

\noindent (in words: the $n$th functional iterate of $\varphi$
evaluated at $z=0$.)
\end{lem}

\begin{proof}
For $n=0$ the claim is true, with the natural
convention that $\varphi^{(0)}(z)=z$. For general $n$,
we have that

\begin{align*}
\mathrm{P}(Z_n = 0) &= \mathbf{E}\left[ \mathrm{P}\left( Z_n=0\,|\,Z_1 \right) \right] = \sum_{k=0}^\infty \mathrm{P}(Z_1=k) \mathrm{P}\left(Z_n=0\,|\,Z_1=k\right) \\
&= \sum_{k=0}^\infty p_k \mathrm{P}\left(Z_n=0\,|\,Z_1=k\right).
\end{align*}

\noindent Now note that
$\mathrm{P}\left(Z_n=0\,|\,Z_1=k\right)=\mathrm{P}(Z_{n-1}=0)^k$,
since that is the probability for $k$ initial
patriarchial specimens to become extinct within $n-1$
generations. So, by induction we get that

\[
\mathrm{P}(Z_n=0) = \sum_{k=0}^\infty p_k \mathrm{P}(Z_{n-1}=0)^k = \varphi(\mathrm{P}(Z_{n-1}=0)) = \varphi(\varphi^{(n-1)}(0))
= \varphi^{(n)}(0).
\]

\end{proof}

\noindent Using the lemma we see that if we are able to prove
that

\[
\mathrm{P}\left(\bigcup_{n=1}^\infty \{Z_n=0\}\right)= \lim_{n\to\infty} \mathrm{P}(Z_n=0) = \lim_{n\to\infty} \varphi^{(n)}(0)<1,
\]

\noindent the final part of Theorem \ref{thm-galton-watson} will
follow. The question therefore reduces to studying the
functional iterates $\varphi^{(n)}(0)$. Note that
$\varphi$ is continuously differentiable on $(0,1)$
and satisfies

\begin{align*}
\varphi(0) &= p_0 \ge 0, \\
\varphi(1) &=1, \\
\varphi'(z) &= \sum_{k=1}^\infty k p_k z^{k-1} \ge 0 \qquad\quad\  \ \textrm{(so $\varphi$ is nondecreasing),}\\
\varphi''(z) &= \sum_{k=2}^\infty k(k-1) p_k z^{k-2} \ge 0 \ \ \textrm{(so $\varphi$ is convex)}, \\
\lim_{z\uparrow 1} \varphi'(z) &= \sum_{k=1}^\infty k p_k = \mu > 1.
\end{align*}

\noindent In fact, since $\mu>1$ there must be some $k\ge 2$
such that $p_k>0$, so $\varphi$ is strictly
increasing and strictly convex.

\begin{xexercise}
Show that from the properties above it follows that
$\varphi$ has a unique fixed point $0\le \rho<1$,
i.e., a point where $\varphi(\rho)=\rho$ (drawing a
rough sketch of the graph of $\varphi$ is very
helpful to explain why this is true).
\end{xexercise}

\begin{proof}[Proof of part 3 of Theorem \ref{thm-galton-watson}]
The sequence $a_n = \varphi^{(n)}(0)$ is increasing,
so converges to a limit $L$. Since $\varphi$ is
increasing, by induction $a_n\le \rho$ for all $n$.
So also
$L=\lim_{n\to\infty} a_n = \mathrm{P}(E) \le \rho <1$,
which was the claim. (In fact, it is easy to see that
$L=\rho$ exactly, since we have

\[
\varphi(L)=\varphi(\lim_{n\to\infty} a_n) = \lim_{n\to\infty} \varphi(a_n) = \lim_{n\to\infty} a_{n+1} = L,
\]

\noindent so $L$ is a fixed point of $\varphi$.)
\end{proof}

\section{The Black-Scholes formula}
In the financial markets, an $\textbf{option}$ is a
contract giving its holder the right to buy or sell an
asset from the issuer of the option for a specified
price at some future date.\footnote{This section is based on sections 15.1--15.2 in the
book \emph{Probability with Martingales} by David Williams.}
Options are useful financial instruments that allow
market participants to calibrate the amount of risk
they want to be exposed to; effectively, risk-averse
players can use options to insure themselves against
unexpected events, paying a premium to other investors
willing to bear the risk (either because they have a
higher tolerance for risk or because their perception
of the amount of risk is different). In theory, this
makes the financial markets more efficient and helps
the economy.

Options have been used in one form or another since
ancient times. However, a good quantitative model
allowing the computation of the monetary value of an
option was not known until two economists, Fischer
Black and Myron Scholes, proposed such a model in 1973
using the mathematics of Brownian motion. Their
formula for the valuation of options, known as the
$\textbf{Black-Scholes formula}$, was the basis for
awarding the 1997 Nobel prize in economics to Scholes
and Robert Merton, another economist who played a part
in the development of the Black-Scholes model (Black
died in 1995 so the prize could not be awarded to
him).

We will discuss here a simplified variant of the
Black-Scholes option pricing model in which time
progresses in discrete steps, so that we do not need
to know about continuous-time processes such as
Brownian motion. Assume that in our simplified
economy, an investor has the option of investing her
money in one of two ways. The first way is to lend the
money; technically, she does this by buying
$\textbf{risk-free bonds}$ (issued by a reliable
entity such as the U.S. government), which are loan
contracts guaranteeing a fixed interest rate of $r$.
Thus, $V$ units of currency invested in bonds at time
$0$ will be worth $(1+r)^n V$ after $n$ time units.
(Of course, this is the result of \emph{compounded
interest}: in this formula we are assuming that after
each time unit our investor takes the interest payment
from the last time unit and reinvests it by buying
more bonds.)

The second type of investment opportunity is to buy
$\textbf{stocks}$. These are assumed to be risky
investments whose value fluctuates randomly. We make
the assumption that the price of stocks is a
multiplicative random walk with an initial (possibly
random) value $S_0$ such that

\[
S_n = (1+R_n) S_{n-1} \qquad (n\ge1),
\]

\noindent where $R_n$ represents the ``random rate of
interest'' in the $n$th time period. We further
assume that $R_1,R_2,\ldots$ are an i.i.d.\ sequence
of random variables satisfying

\[
\mathrm{P}(R_n = a) = 1-p, \ \ \mathrm{P}(R_n = b) = p,
\]

\noindent where $a,b$ are two values in $(-1,\infty)$
satisfying $a<r<b$, and $p$ is taken to be
$p=\frac{r-a}{b-a}$. This choice of $p$ ensures that
$\mathbf{E}(R_1) = r$, so that on average the random
interest rate is equal to the risk-free interest rate.
(Any other choice of $p$ would cause an investor to
always prefer one of the two asset classes to the
other, making the model essentially uninteresting.)

Our investor starts with an amount $x$ of money and
invests it over time, allocating the available funds
to either stocks or bonds as she wishes. Formally, let
$\mathcal{G}_n = \sigma(S_0,R_1,\ldots,R_n)$. Denote
$B_0=1, B_n=(1+r)^n$ (in analogy with the
multiplicative random walk $S_n$, this is the value
of an investment account concentrated in bonds and
started from an initial investment of 1 unit of
currency). An $\textbf{investment strategy}$ is a
pair $(A_n, V_n)_{n=0}^N$ of processes (where $N$ may
be finite, or infinite) which are predictable with
respect to the filtration $(\mathcal{G}_n)_n$, and
such that the following relations are satisfied:

\begin{align*}
x &= \tag{9}\label{eq:investment-strategy-initial} A_0 S_0 + V_0 B_0, \\
A_n S_n + V_n B_n &= \tag{10}\label{eq:investment-strategy-rebalance} A_{n+1} S_n + V_{n+1} B_n, \qquad (1\le n< N).
\end{align*}

\noindent Here, $A_n$ denotes the number of stock investment
units held between time $n$ and $n+1$, and $V_n$
denotes the number of bond investment units held
between time $n$ and $n+1$. Denote
$X_n = A_n S_n + V_n B_n$. This represents the value
of the investment account at time $n$. The relation
(\ref{eq:investment-strategy-initial}) means that the
initial value $X_0$ of the investment account is $x$
units of currency, and
(\ref{eq:investment-strategy-rebalance}) represents the
fact that immediately after time $n$ the investor
rebalances the investment portfolio, changing the
mixture of bonds and stocks but keeping the same
amount invested. In a slight simplification of real
life, the model assumes that the act of rebalancing
the portfolio does not involve paying a fee to a
brokerage or account management company, an assumption
technically referred to as
$\textbf{zero transaction costs}$.

In our model, we now add a third type of investment
opportunity. A $\textbf{European option}$ is a
contract that gives an investor the right (but not the
obligation) to buy $1$ unit of stocks at some fixed
time $N$ for a price of $K$ units of currency. The
act of taking advantage of this right is referred to
as $\textbf{exercising}$ the
option\footnote{A common variant is the $\textbf{American option}$,
in which the right to exercise the option exists at
any time up to and including time $N$.}. The time $N$ is called
the $\textbf{expiration time}$, and $K$ is called the
$\textbf{strike price}$. The question we wish to
address is: what is the ``fair value''
$v_\textrm{fair}$ such that an investor would
consider paying an amount $v_\textrm{fair}$ (or
better yet, slightly less than $v_\textrm{fair}$) for
such an option offered to him at time $0$? Note that
at time $N$ the option is worth $S_N-K$ if $S_n\ge K$,
or $0$ otherwise (i.e., if the market price of stocks
at time $N$ is below the strike price at expiration
time, the option gives a pointless right to buy stocks
at a higher price than the one for which they are
being offered for sale on the free market. This is a
right which no rational investor would want to
exercise; in this case the option is said to have
$\textbf{expired worthless}$). In other words, the
value of the option at its expiration time is exactly
the (random) amount $(S_N-K)_+$, the positive part of
the random variable $S_N-K$. An investor considering
an investment in options at time $0$ would be
reasonable to consider the expected value of this
quantity,

\[
\mathbf{E}(S_N-K)_+,
\]

\noindent as her expected return after $N$ time units of
investment. Comparing this to the return on investing
the same amount $v_\textrm{fair}$ in risk-free bonds,
which is equal to

\[
(1+r)^N v_\textrm{fair},
\]

\noindent the investor may conclude that $v_\textrm{fair}$ is
exactly the value that equalizes both quantities (thus
making her indifferent to the choice of which type of
investment to make), namely

\begin{equation*} \tag{11}\label{eq:discrete-black-scholes}
v_\textrm{fair} = (1+r)^{-N} \mathbf{E}(S_N-K)_+.
\end{equation*}

\noindent In finance, it is typical to think of this type of
expression as the quantity $\mathbf{E}(S_N-K)_+$
``discounted'' $N$ units of time into the future. The
idea is that the promise of future money is worth less
than actual money in the present. The ``rate of
discounting,'' which is the factor by which it becomes
worth less and less with every time unit it recedes
further into the future, is exactly the risk-free
interest rate $1+r$. (Think of a situation in which
someone offers you a choice between \(\$100\) today
or \(\$150\) in a month's time. Which would you
prefer? What would be your criterion for deciding if
the number $150$ were replaced with any other number
higher than $100$?)

Equation (\ref{eq:discrete-black-scholes}) is our
discrete version of the Black-Scholes formula, but the
reasoning that lead us to it is a bit shaky, and can
be strengthened in the following way. Assume that a
market for options hasn't developed yet, so none are
being offered for sale. An investor may come up with a
clever way of replicating the outcome of investing in
an option using a particular way of managing a
portfolio of stocks and bonds. Formally, a \textbf{hedging
strategy with initial value $x$} is a pair of
processes $(A_n,V_n)_{n=0}^N$ which is an investment
strategy in the sense defined above, with an initial
value $x$ as in (\ref{eq:investment-strategy-initial}),
and such that in addition we have the properties

\begin{align*}
X_n &\ge 0,  \qquad \qquad \qquad (0\le n\le N), \\
X_N &= (S_N-K)_+,
\end{align*}

\noindent (where as before $X_n = A_n S_n + V_n B_n$ denotes
the value of the portfolio at time $n$). Note that a
hedging strategy behaves for all intents and purposes
like an option with strike price $K$ and expiration
time $N$; any investor who possesses such a strategy
could issue actual options with these parameters and
sell them to other investors, charging slightly more
than $x$ units of currency, investing $x$ of them in
the hedging scheme (which precisely offsets the
options he sold) and making an instant and risk-free
profit. (Note that the condition $X_n\ge0$ is
required to ensure that the manager of the hedging
strategy will not have to inject new cash at any point
in time to keep the investment going; she will never
be ``underwater.'') Thus, we have a rather convincing
argument that any value of $x$ for which one is able
to find such a hedging strategy is, \emph{by definition},
the fair value of the option.

\begin{thm}
A hedging strategy with initial value $x$ exists if
and only if

\begin{equation*} \tag{12}\label{eq:hedging-strategy-condition}
x=(1+r)^{-N}\mathbf{E}(S_N-K)_+,
\end{equation*}

\noindent and in this case it is unique.
\end{thm}

\begin{proof}
Assume that there is a hedging strategy with initial
value $x$. Let $Y_n = (1+r)^{-n} X_n$ (the discounted
value of $X_n$ at time $0$). The proof will depend on
showing that $(Y_n)_n$ is a martingale. To see this,
note that from (\ref{eq:investment-strategy-rebalance}) we
have that

\begin{align*}
X_n - X_{n-1} &= (A_n S_n + V_n B_n) - (A_n S_{n-1} + V_n B_{n-1})
= A_n (S_n-S_{n-1}) + V_n (B_n-B_{n-1}) \\
&= A_n R_n S_{n-1} + r V_n B_{n-1} = A_n S_{n-1}(R_n-r) + r(A_n S_{n-1} + V_n B_{n-1}) \\
&= A_n S_{n-1} (R_n-r) + r X_{n-1}.
\end{align*}

\noindent It follows that

\[
Y_n-Y_{n-1} = (1+r)^{-n}(X_n-(1+r)X_{n-1}) = (1+r)^{-n}A_n S_{n-1}(R_n-r).
\]

\noindent Thus, $Y_n$ can be expressed as

\begin{align*}
Y_n &= Y_0 + \sum_{k=1}^n (Y_k-Y_{k-1}) =
x + \sum_{k=1}^n (1+r)^{-k}A_k S_{k-1} (R_k-r) \\
&= x + \sum_{k=1}^n F_k (Z_k-Z_{k-1}),
\end{align*}

\noindent where $F_n = (1+r)^{-n} A_n S_{n-1}$ and
$Z_n = \sum_{k=1}^n (R_k-r)$. In other words, we have
shown that $Y_n$ can be represented as

\[
Y_n = x + (F\bullet Z)_n.
\]

\noindent Here, the sequence $(F_n)_n$ is predictable with
respect to the filtration $(\mathcal{G}_n)_n$, and
$Z_n$ is a martingale, so this is a genuine
martingale transform, and $(Y_n)_n$ is a martingale,
as claimed above. In particular we get (using the
definition of a hedging strategy) that

\[
x = Y_0 = \mathbf{E}(Y_0) = \mathbf{E}(Y_N) = (1+r)^{-N} \mathbf{E}(X_N) = (1+r)^{-N} \mathbf{E}(S_N-K)_+,
\]

\noindent proving (\ref{eq:hedging-strategy-condition}). This
completes the ``only if'' part of the proof.

Conversely, assume that
$x=(1+r)^{-N} \mathbf{E}(S_N-K)_+$. We need to
construct a hedging strategy with this initial value,
and show that it is uniquely determined. Inspired by
the insights gained by the computation above, we
define a martingale

\[
Y_n = \mathbf{E}\left((1+r)^{-N}(S_N-K)_+\,|\,\mathcal{G}_n\right), \qquad (0\le n\le N)
\]

\noindent so that $Y_0 = x$ and $Y_N = (1+r)^{-N} (S_N-K)_+$,
and look for a predictable sequence $(F_n)_n$ such
that $Y_n = x + (F\bullet Z)_n$. The fact that a
unique such process exists is the result of a simple
computation which can be distilled into the following
lemma.

\begin{lem}
\label{lem-easy-martingale-rep}
If $(M_n)_{n=0}^N$ is a martingale w.r.t.\ the
filtration $(\mathcal{G}_n)_n$ satisfying $M_0 = 0$,
then there is a unique predictable process
$(F_n)_{n=1}^N$ such that $M = F\bullet Z$, with
$Z_n=\sum_{k=1}^n (R_k-r)$ as above.
\end{lem}

\begin{xexercise}
Prove Lemma \ref{lem-easy-martingale-rep}.
\end{xexercise}

\noindent Working our way backwards, we now define
$X_n = (1+r)^n Y_n$, $A_n = (1+r)^n F_n/S_{n-1}$ and
$V_n = (X_{n-1}-A_n S_{n-1})/B_{n-1}$. Then
$X_n\ge 0$, $(A_n)_n$ and $(V_n)_n$ are predictable
sequences, and the above computations can be read in
reverse to show that
(\ref{eq:investment-strategy-rebalance}) is satisfied,
i.e., the sequences $(A_n, V_n)_n$ represent a valid
hedging strategy.
\end{proof}

\noindent $\textbf{Notes.}$

\begin{enumerate}[label=\arabic*.]
\item In the proof above, one can check that the processes
$A_n, V_n$ are nonnegative. This means that, although
we didn't require it as part of the definition of
hedging strategies, in practice hedging will never
involve investing negative amounts of money (which
would correspond to ``shorting'' stocks or issuing
bonds).

\item The formula (\ref{eq:discrete-black-scholes}) is written
in a form that is not very explicit, as it requires
computing a messy average for a multiplicative random
walk. In the ``real'' Black-Scholes formula, where the
stock price $(S_n)_n$ process is modeled by a
$\textbf{geometric Brownian motion}$ (the
continuous-time analogue of a multiplicative random
walk), one can evaluate the corresponding expectation
to obtain a much more explicit formula. The
Black-Scholes model has also been adapted to
American-style options and other variants. As usual,
if you want to learn more about this subject,
Wikipedia will make you wish you didn't.
\end{enumerate}

